\scanpage{222}
linéaire sur $\mathcal M$, on trouve une unique antidérivation $\delta$ de $\Lambda^*\mathcal M$ qui sur $\Lambda^1\mathcal M$ se réduise à $\varphi$, et on a $\delta^2=0$. (N.B. On a $H_0(\Lambda^*\mathcal M,\delta)=S/\varphi(\mathcal M)$, et si $\mathcal M$ est libre de t.f., alors $H_i(\Lambda^*\mathcal M,\delta)=0$ pour $i\neq0$ ssi $\mathcal M\longrightarrow S$ est ``transverse'' -- il suffit p.c. qu'une base \unsure{choisie} de $\mathcal M$ définisse une suite $S$-régulière, et c'est aussi néc. si $S$ noeth et le $\varphi(\mathcal M)\subset\operatorname{rad}S$)

Ceci dit, on obtient les complexes $(2)$ et $(4)$ en prenant respectivement $S=\operatorname{Sym}^*M$ ou $S=\Gamma^*M$, $\mathcal M=M\otimes_A S$, et en prenant l'hom $\mathcal M\longrightarrow S$ \deleted{\unsure{le}} dont la restriction à $M$ est l'hom can $M\xrightarrow{\sim}\deleted{\operatorname{Sym}^1M\subset S}\inserted{S^1\subset S}$ \deleted{$\unsure{\operatorname{Sym}^1M}$} $(S^1=\operatorname{Sym}^1M\text{ resp }\Gamma^1M)$.

Les opérateurs $d\delta\unsure{,|+}\delta d$ dans $\operatorname{Sym}^*\otimes\Lambda^*$ resp $\Gamma^*\otimes\Lambda^*$ ne sont pas bien \unsure{jolis}, mais on voit de suite (tout \unsure{anticommut} des antidérivations est une dérivation) que
\[
\theta=d\delta+\delta d
\]
est une dérivation, \deleted{et} \deleted{\unsure{pour le cas $S^*=\Gamma^*$ est compatible aux p.d.}} ne changeant \deleted{pas le} le bidegré, et qui sur \deleted{$\unsure{\Lambda^1}$} $S^1=\operatorname{Sym}^1$ ou $\Gamma^1$ et sur $\Lambda^1$ est l'identité; \inserted{de plus si $S^*=\Gamma^*$, il est compatible aux puiss. div. sur $\Gamma^*$}. Donc on en conclut facilement que l'on a
\[
\boxed{(d\delta+\delta d\unsure{)}x=\theta x=nx}
\]
\margin{$x$ de degré \uline{total} $n$ i.e. $x\in\sum_{\substack{p+q=n\\p,q\geq0}}S^p\otimes\Lambda^q$ \unsure{(hom}}
(Cela montre donc que si $A$ est de car. \unsure{$n$ inversible $\forall n>0$}) alors $H^n_{DR}(S^*\otimes\Lambda^*)$ et $H_n^{\mathrm{Kos}}(S^*\otimes\Lambda^*)$ ($S^*$ étant $\operatorname{Sym}^*\simeq\Gamma^*$) est nul si \deleted{\unsure{donc}} $n>0$ (sans hyp. de platitude!). Plus gén., $H^n_{DR}$ et $H_n^{\mathrm{Kos}}$ sont nuls pour tout $n$ tel que $n.1_A$ inv. dans $A$.
\figure{\begin{tikzpicture}\draw (0,0)--(0,1);\draw (.08,0)--(.08,1);\end{tikzpicture}}
\margin{221}
